% A minimal thesis skeleton using the standard `report` class.
% Exercises chapter / section structure, theorem environments, and a TOC.
\documentclass[12pt,a4paper]{report}

\usepackage[utf8]{inputenc}
\usepackage[T1]{fontenc}
\usepackage{geometry}
\usepackage{amsmath}
\usepackage{amssymb}
\usepackage{amsthm}
\usepackage{graphicx}
\usepackage{hyperref}
\usepackage{natbib}

\newtheorem{theorem}{Theorem}[chapter]
\newtheorem{lemma}[theorem]{Lemma}
\newtheorem{corollary}[theorem]{Corollary}
\newtheorem{definition}[theorem]{Definition}

\title{A Generic Thesis Skeleton}
\author{Alice Anderson}
\date{June 2026}

\begin{document}

\maketitle

\tableofcontents

\listoffigures

\chapter{Introduction}

This thesis presents a small skeleton designed to exercise ByeTex on report-class
documents. We use chapters as the top-level structural unit, with sections and
subsections nested below.

\section{Motivation}

Mathematical typesetting in LaTeX is a long-standing standard; Typst is the new
contender. Migrating a thesis between the two should not require manual rewrite.

\section{Contributions}

\begin{itemize}
\item A summary of the migration challenges, including macro expansion and table layout.
\item An evaluation of conversion fidelity across three baseline templates.
\item A discussion of the residual long tail and where manual intervention is unavoidable.
\end{itemize}

\chapter{Background}

\section{Notation}

Throughout this thesis, we use the following conventions. Sets are denoted by
capital letters such as $A$, $B$, $C$. The cardinality of a set $A$ is $|A|$.

\begin{definition}[Metric space]
A metric space is a pair $(X, d)$ where $X$ is a set and $d : X \times X \to
\mathbb{R}_{\geq 0}$ satisfies the usual three axioms.
\end{definition}

\begin{theorem}[Cauchy criterion]
\label{thm:cauchy}
A sequence $(x_n)$ in a complete metric space converges if and only if it is
Cauchy.
\end{theorem}

\begin{proof}
The forward direction is immediate from the triangle inequality. The reverse
direction follows from completeness; see~\cite{rudin1976}.
\end{proof}

\chapter{Method}

\section{Setup}

By Theorem~\ref{thm:cauchy}, any Cauchy sequence converges. We exploit this
fact to define the limit
\begin{equation}
x^* = \lim_{n \to \infty} x_n.
\label{eq:limit}
\end{equation}

\section{Algorithm}

The procedure iterates equation~\eqref{eq:limit} until the residual falls below
$10^{-6}$. In practice, this converges within $200$ iterations on the test
problems considered in this thesis.

\chapter{Results}

We benchmark against the three baselines from \cite{einstein,dirac,bohr}.
Performance is summarized in the next chapter.

\chapter{Conclusion}

This skeleton intentionally avoids exotic class files. The structure should
convert cleanly through ByeTex with the standard mappings for chapters,
sections, theorems, equations, and references.

\bibliographystyle{plainnat}
\bibliography{thesis_refs}

\end{document}
